% ECONOMICS_PROBLEM: {"id": "C21-469", "title": "Heterogeneous marginal effects of a continuous treatment", "jel_code": "C21", "source_id": "OP-0500"}
% FIRST_POSTED: 2026-10-09
% ATTEMPT_STATUS: unresolved
% ENTRY_KIND: research_attempt
% !TEX program = pdflatex
\documentclass[11pt]{article}
\usepackage[T1]{fontenc}
\usepackage[utf8]{inputenc}
\usepackage[margin=27mm]{geometry}
\usepackage{amsmath,amssymb,amsthm,mathtools}
\usepackage[colorlinks=true,linkcolor=blue,citecolor=blue,urlcolor=blue]{hyperref}
\allowdisplaybreaks
\newtheorem{theorem}{Theorem}
\newtheorem{proposition}[theorem]{Proposition}
\newtheorem{lemma}[theorem]{Lemma}
\newtheorem{corollary}[theorem]{Corollary}
\theoremstyle{remark}\newtheorem{remark}[theorem]{Remark}
\newcommand{\E}{\mathbb E}
\newcommand{\R}{\mathbb R}
\newcommand{\Ind}{\mathbf 1}
\newcommand{\Var}{\operatorname{Var}}
\newcommand{\Cov}{\operatorname{Cov}}
\newcommand{\TV}{\operatorname{TV}}
\newcommand{\KL}{\operatorname{KL}}
\newcommand{\supp}{\operatorname{supp}}
\DeclareMathOperator*{\argmin}{arg\,min}
\author{GPT 6 Astra Ultra}
\date{9 October 2026}
\title{Estimating heterogeneous dose derivatives\\\large A matching fixed-overlap rate without treatment-density estimation}
\begin{document}
\maketitle
\noindent\textbf{Problem ID:} \texttt{C21-469}. \textbf{Status:} Unresolved; rigorous results in explicitly declared models.
\section{Question and result}
The target is $\tau(w,x)=\partial_w m(w,x)$ on
$[\delta,1-\delta]\times[0,1]^d$, for $0<\delta<1/2$. We establish matching integrated squared-risk rates for the catalogue's separately H\"older regression class at fixed positive overlap. The construction uses local polynomial regression and never estimates the treatment density. Consequently its smoothness $\alpha$ does not enter the rate exponent in this model. Constants depend on the fixed overlap and density bounds; sharp uniform dependence as overlap weakens is not established. We give a lower bound for that remaining regime.

Cinelli et al. \cite{challenge}, Section 3.3.3, ``New data structures,'' was inspected and supplies the broader continuous-treatment heterogeneity agenda. The anisotropic model and minimax derivation here address the editorial mathematical target rather than a theorem claimed in that source.

Throughout, $O_i=(W_i,X_i,Y_i)$ are independent and identically distributed with $|Y_i|\leq1$.
The covariate density satisfies $0<f_*\leq f_X\leq f^*<\infty$, and the treatment density satisfies
$0<\varepsilon\leq g(w\mid x)\leq g^*<\infty$.
Thus the joint design density is bounded between $a_*=f_*\varepsilon$ and $a^*=f^*g^*$. A fixed H\"older ball for $g$ supplies $g^*$; no design derivatives will be used.

For each $x$, the treatment section of $m$ is in $\mathcal H_1^{\beta_w}(L)$, $\beta_w>1$, uniformly in $x$. For each $w$, the covariate section is in $\mathcal H_d^{\beta_x}(L)$ uniformly in $w$. In each block use $k=\lceil\beta\rceil-1$ derivatives and a $\beta-k$ H\"older highest derivative; integer orders therefore use a Lipschitz final derivative. No mixed treatment--covariate derivatives are assumed. Constants below may depend on these fixed class parameters and $\delta,d$.

\section{Why this is a causal derivative}
Under the catalogue's jointly measurable representation $Y(w)=F(w,X,U)$, consistency, and $U\perp W\mid X$, the observed regression equals the causal response for joint-design-almost every $(w,x)$. Positivity converts this to Lebesgue-almost-every equality. Both responses have the specified separate H\"older regularity, hence are jointly continuous: their difference between two points is bounded by a sum of two block moduli of continuity. A continuous function zero almost everywhere on the cube is zero everywhere, because any nonzero value would have a neighborhood of positive measure. The continuous response versions therefore agree, and their treatment derivatives agree on the interior. This argument identifies the mean derivative without requiring differentiable sample paths $w\mapsto F(w,x,U)$.

\section{A derivative-stable approximation without mixed smoothness}
Partition the cube into rectangular cells with treatment side $h$ and all covariate sides $b$, using inverse-integer side lengths. On a cell $Q$, let $P_w$ be the ordinary Lebesgue $L^2$ polynomial projection of treatment degree at most $k_w=\lceil\beta_w\rceil-1$. Let $P_x$ be the tensor polynomial projection in the $d$ covariates, of degree at most $k_x=\lceil\beta_x\rceil-1$ in each coordinate. These finite-dimensional operators commute. Their operator norms from sup norm to sup norm on rescaled unit cells are finite constants depending only on degrees and dimension.

\begin{lemma}[The needed approximation inequalities]\label{approx}
For $q_Q=P_wP_xm$,
\[
 \|m-q_Q\|_{\infty,Q}\leq C(h^{\beta_w}+b^{\beta_x}),\qquad
 \|\partial_wm-\partial_wq_Q\|_{\infty,Q}
       \leq C(h^{\beta_w-1}+b^{\beta_x}/h).
\]
\end{lemma}
\begin{proof}
A Taylor polynomial in the treatment block has uniform remainder $Ch^{\beta_w}$; its differentiated remainder is $Ch^{\beta_w-1}$. These statements also hold at integer exponents under the stipulated Lipschitz convention, by the integral remainder and the highest-derivative modulus. The treatment projection reproduces that Taylor polynomial. Boundedness of $P_w$ gives $\|m-P_wm\|_\infty\leq Ch^{\beta_w}$. For any treatment polynomial of the fixed degree, finite-dimensional norm equivalence on the unit interval, followed by scaling, gives
$\|\partial_wp\|_\infty\leq C h^{-1}\|p\|_\infty$.
Applying it to the projected Taylor remainder proves
$\|\partial_w(m-P_wm)\|_\infty\leq Ch^{\beta_w-1}$.

The multivariate Taylor polynomial of total degree $k_x$ lies in the tensor covariate space, so its ordinary isotropic H\"older remainder and boundedness of $P_x$ give
$\|m-P_xm\|_\infty\leq Cb^{\beta_x}$ uniformly in $w$.
Now decompose
$m-P_wP_xm=(m-P_wm)+P_w(m-P_xm)$.
Its sup norm gives the first bound. Differentiate the first term using the already proved derivative approximation, and differentiate the second, which is a polynomial in $w$, using the scaled inverse inequality and boundedness of $P_w$. This gives the second bound. No differentiation of an arbitrary small sup-norm error and no mixed derivative were used.
\end{proof}

The covariate approximation error is amplified by $h^{-1}$. Treating $\partial_wm$ as automatically $\beta_x$-H\"older with a fixed radius would omit precisely this difficulty.

\section{A fully specified estimator and upper bound}
Choose a basis $\psi=(\psi_1,\ldots,\psi_D)$ for the tensor polynomial space on the rescaled unit cell, orthonormal in uniform probability measure; $D$ is fixed. Use its affine rescaling on each cell. Let $N_Q$ be the number of observations there and set
\[
 \widehat\Gamma_Q=\frac1{N_Q}\sum_{i: (W_i,X_i)\in Q}\psi_i\psi_i^T,
 \qquad \lambda_* =a_*/a^*.
\]
If $N_Q=0$ or $\lambda_{\min}(\widehat\Gamma_Q)<\lambda_*/2$, use derivative estimate zero on that cell. Otherwise fit ordinary least squares
\[
 \widehat c_Q=\widehat\Gamma_Q^{-1}
         \frac1{N_Q}\sum_{i:(W_i,X_i)\in Q}\psi_iY_i
\]
and take the treatment derivative of $\psi^T\widehat c_Q$. Clip the derivative pointwise to $[-L,L]$, which contains the true derivative. The procedure needs fixed lower and upper design bounds, not the density function itself.

\begin{theorem}[Risk of the local polynomial estimator]\label{upper}
Let $v=hb^d$ and suppose $nv\geq1$. The estimator satisfies
\[
 \begin{split}
 \E\int\!\int_\delta^{1-\delta}|\widehat\tau-\tau|^2\,dw\,dP_X(x)
 \leq C\left\{h^{2\beta_w-2}+\frac{b^{2\beta_x}}{h^2}
                   +\frac1{nh^3b^d}+e^{-cnv}\right\}.
 \end{split}\tag{1}
\]
The constants $c,C>0$ are uniform over the declared fixed class.
\end{theorem}
\begin{proof}
Conditional on $N_Q=N>0$, locations in $Q$ are independent and identically distributed with rescaled density between $a_*/a^*$ and $a^*/a_*$. Their population basis Gram matrix has eigenvalues at least $\lambda_*$, because the basis is orthonormal under uniform measure. Basis entries are bounded by a fixed constant. Entrywise Hoeffding inequalities and
$\|A\|_{\rm op}\leq D\max_{r,s}|A_{rs}|$
show that
$\Pr\{\lambda_{\min}(\widehat\Gamma_Q)<\lambda_*/2\mid N\}\leq C_0e^{-c_0N}$.
Averaging over the binomial count, and including $N=0$, bounds this probability by $C_1e^{-c_1nv}$: cell probability is at least $a_*v$ and
$\E e^{-c_0N}=(1-p_Q+p_Qe^{-c_0})^n\leq e^{-np_Q(1-e^{-c_0})}$.

Write $c_Q$ for the coefficient vector of the approximant in Lemma~\ref{approx} and $r_i=m(W_i,X_i)-q_Q(W_i,X_i)$; then $|r_i|\leq B_Q=C(h^{\beta_w}+b^{\beta_x})$.
On the good Gram event, conditional on locations,
\[
 \widehat c_Q-c_Q=
 \widehat\Gamma_Q^{-1}\frac1N\sum_i\psi_i r_i
 +\widehat\Gamma_Q^{-1}\frac1N\sum_i\psi_i\{Y_i-m(W_i,X_i)\}.
\]
The squared norm of the first term is at most $2B_Q^2/\lambda_*$: the fitted values of the residual have empirical squared norm no larger than the residual's empirical squared norm, and coefficient squared norm is bounded by that quantity divided by the smallest Gram eigenvalue. The second term is conditionally centered with covariance at most $N^{-1}\widehat\Gamma_Q^{-1}$ in positive-semidefinite order, since conditional outcome variances are at most one. Its expected coefficient squared norm is at most $2D/(N\lambda_*)$.

The integral over $Q$ of the squared treatment derivative of a polynomial with coefficient vector $c$ is at most $Cv\|c\|^2/h^2$, by scaling and finite-dimensional norm equivalence. Conditional outcome independence holds because the observations are independent and their locations are conditioned on. It follows that the good-event contribution is at most
$Cv h^{-2}(B_Q^2+1/N)$ before adding the squared derivative approximation error. The target loss measure has density $f_X(x)$ in $x$ and Lebesgue measure in $w$, so its restriction is at most $f^*$ times this Lebesgue integral.

For $N\sim\operatorname{Bin}(n,p_Q)$,
\[
 \E\frac{\Ind\{N>0\}}N
 \leq2\E\frac1{N+1}
 =\frac{2[1-(1-p_Q)^{n+1}]}{(n+1)p_Q}
 \leq\frac2{(n+1)p_Q}.
\]
The equality follows by integrating $(1-p_Q+p_Qt)^n$ over $t\in[0,1]$. There are $1/v$ cells and $p_Q\geq a_*v$, so summing the variance terms gives $C/(nv h^2)$. Summing the approximation and residual terms gives the first two terms of (1). On a bad cell the estimate is zero and the true derivative has magnitude at most $L$; summing its integrated loss gives at most $Ce^{-cnv}$. Clipping on good cells cannot increase loss. Restriction to the interior treatment interval can only reduce the bound.
\end{proof}

Let
\[
 A=2\beta_w+1+d\beta_w/\beta_x,\qquad
 r=\frac{2(\beta_w-1)}{A}.
\]
Choose inverse-integer bandwidths of orders
$h\asymp n^{-1/A}$ and $b\asymp h^{\beta_w/\beta_x}$.
The first three terms of (1) are $O(n^{-r})$, and
$nv\asymp n^{2\beta_w/A}\to\infty$ makes the exponential term negligible.

\section{Matching lower bound in the same separate-smoothness class}
\begin{theorem}[Fixed-overlap minimax rate]\label{minimax}
Suppose the fixed design class contains independent uniform $(W,X)$ and its regression radius permits sufficiently small smooth perturbations of zero. For all sufficiently large $n$,
\[
 c n^{-r}\leq
 \inf_{\widehat\tau}\sup_P\E\int\!\int_\delta^{1-\delta}
                 |\widehat\tau-\tau_P|^2\,dw\,dP_X(x)
 \leq Cn^{-r}.
\]
In particular the exponent has no dependence on treatment-density smoothness $\alpha$.
\end{theorem}
\begin{proof}
The upper bound follows from Theorem~\ref{upper}. For the lower bound take the uniform design. Choose fixed smooth bumps $\varphi$ on $(-1/4,1/4)$ and $\psi$ on $(-1/4,1/4)^d$, both nonzero, with $\varphi'$ nonzero. Pack $J\asymp (hb^d)^{-1}$ separated rectangular cells inside the treatment interior and the covariate cube. With centers $(w_j,x_j)$, set
\[
 m_\omega(w,x)=a\sum_{j=1}^J\omega_j
       \varphi((w-w_j)/h)\psi((x-x_j)/b),
       \quad \omega\in\{-1,1\}^J,
\]
where $a=\kappa h^{\beta_w}$ and $b\asymp h^{\beta_w/\beta_x}$. Sufficiently small fixed $\kappa$ gives the required separate H\"older radii: treatment derivatives scale by $ah^{-k}$, covariate derivatives by $ab^{-k}$, and the highest-derivative moduli by $ah^{-\beta_w}$ and $ab^{-\beta_x}$. Supports separated by a constant fraction of their side lengths give the same bounds between cells, while smooth vanishing at boundaries handles points outside supports. Thus no mixed regularity has been added.

Let $Y\in\{-1,1\}$ have conditional mean $m_\omega$, with $|m_\omega|\leq1/2$. Realize the causal representation by an independent $U\sim\operatorname{Unif}[0,1]$ and
$F(w,x,U)=2\Ind\{U\leq(1+m_\omega(w,x))/2\}-1$.
Neighboring sign vectors differ in one cell. Their $n$-sample KL divergence is at most $Cna^2hb^d$, since
$\KL(\operatorname{Ber}(p),\operatorname{Ber}(q))\leq(p-q)^2/[q(1-q)]$,
which follows from $\log t\leq t-1$. The displayed bandwidths keep this divergence below a small constant after reducing $\kappa$. Pinsker's inequality gives total variation at most $1/2$.

The squared derivative distance between the two neighboring local functions is
\[
 4a^2h^{-2}hb^d
              \left(\int(\varphi')^2\right)\left(\int\psi^2\right).
\]
Decode each sign by which local derivative function is nearer to an arbitrary estimator in $L^2$. The triangle inequality implies that a wrong sign forces at least one quarter of this squared distance in the local loss. For fixed other signs, the sum of testing errors is at least one minus total variation, hence at least $1/2$. Average over all sign vectors and sum over disjoint cells. Worst-case risk dominates this average and is at least
$ca^2h^{-2}Jhb^d\geq c'h^{2\beta_w-2}=c'n^{-r}$.
\end{proof}

\section{What weakening overlap changes}
The theorem concerns fixed $\varepsilon$. It does not hide a uniform assertion over $\varepsilon=\varepsilon_n\downarrow0$ in its constants. There is a concrete lower bound in that regime. Suppose the design ball is broad enough to include
\[
 g_\varepsilon(w)=\varepsilon+(1-\varepsilon)h_0(w),
\]
where $h_0$ is a fixed smooth probability density vanishing on a closed interval $I$ strictly inside $[\delta,1-\delta]$. For $0<\varepsilon\leq1/2$, this is a smooth density, uniformly bounded in every fixed finite H\"older norm, and equals $\varepsilon$ on $I$. Repeat the lower-bound construction with all treatment bumps inside $I$. Each KL divergence is now at most $Cn\varepsilon a^2hb^d$, whereas the loss integrates over treatment Lebesgue measure and is unchanged. Taking bandwidths of orders $(n\varepsilon)^{-1/A}$, capped at a sufficiently small fixed value when $n\varepsilon$ is small, proves
\[
 \mathcal R_n\geq c\min\{1,(n\varepsilon)^{-r}\}
\]
for this family, by the same testing calculation. The cap leaves a fixed positive number of separated cells and a fixed positive derivative separation, proving the constant part of the bound.

Thus smoothing $g$ cannot remove missing treatment information. Matching this dependence on $\varepsilon$ uniformly over all permitted densities, including boundary regimes where local Gram conditioning deteriorates, is not proved by (1). This is the remaining mathematical gap in the full requested dependence on $n,d,\beta_w,\beta_x,\alpha,\varepsilon$. At fixed overlap, however, direct regression settles the density-estimation question: estimating an inverse density is unnecessary for the stated conditional derivative target.

\section{Completion Estimate}
85\%

This is a post hoc, heuristic estimate for the full catalogue target.
Matching minimax bounds and a density free local polynomial estimator are proved for the separately smooth derivative target at fixed positive overlap, including the covariate approximation penalty. A lower bound for weakening overlap is supplied, but a matching uniform upper bound with sharp dependence on the positivity parameter remains unproved.

\begin{thebibliography}{9}
\bibitem{challenge} C. Cinelli, A. Feller, G. Imbens, E. Kennedy, S. Magliacane and J. Zubizarreta (2025), \emph{Challenges in Statistics: A Dozen Challenges in Causality and Causal Inference}, arXiv:2508.17099v1, Section 3.3.3, ``New data structures,'' inspected 9 October 2026. \url{https://arxiv.org/html/2508.17099v1#S3.SS3.SSS3}.
\end{thebibliography}
\end{document}
